\section{Methodology}
\label{sec:method}

\subsection{Task Formulation}

We frame the problem as a binary classification task: given an item $x$ and a category $c$, does the item exist as a member of category $c$?
For each probe, we query a language model $\gM$ with a prompt $p(x, c)$ and parse the response as a binary judgment (exists / does not exist).
We measure the \emph{false positive rate} (\fpr)---the probability that the model claims a nonexistent item exists---as our primary metric.

Formally, let $\gS^+ = \{x : x \in c\}$ be the set of real members and $\gS^- = \{x : x \notin c\}$ be nonexistent items. We define:
\begin{align}
    \text{FPR} &= P(\gM(p(x,c)) = \text{``exists''} \mid x \in \gS^-) \\
    \text{TPR} &= P(\gM(p(x,c)) = \text{``exists''} \mid x \in \gS^+)
\end{align}

\subsection{Probe Dataset}

We construct a dataset of 162~items across 7~categories, each containing three item types: \emph{real} (ground-truth members), \emph{plausible nonexistent} (real-world entities that could plausibly belong but do not), and \emph{implausible nonexistent} (entities unlikely to belong).
\tabref{tab:dataset} summarizes the dataset composition.

\begin{table}[t]
    \centering
    \caption{Probe dataset composition across 7~categories. Density refers to the number of real members in each category. Ground truth for emoji categories is Unicode Emoji v16.0 \cite{unicode2024emoji}; geography and chemistry use standard references.}
    \begin{tabular}{@{}lcccr@{}}
        \toprule
        \textbf{Category} & \textbf{Real} & \textbf{Plausible} & \textbf{Implausible} & \textbf{Density} \\
        \midrule
        Marine animal emojis & 15 & 10 & 5 & High \\
        Land animal emojis   & 15 & 10 & 5 & High \\
        Bird emojis          & 10 &  8 & 4 & Medium \\
        Insect emojis        &  8 &  6 & 4 & Medium \\
        Fruit emojis         & 14 &  8 & 4 & High \\
        Country capitals     & 10 &  5 & 3 & High \\
        Chemical elements    & 10 &  5 & 3 & High \\
        \midrule
        \textbf{Total}       & \textbf{82} & \textbf{52} & \textbf{28} & --- \\
        \bottomrule
    \end{tabular}
    \label{tab:dataset}
\end{table}

\para{Plausibility criteria.}
``Plausible nonexistent'' items are real-world entities that share strong semantic similarity with existing category members (\eg seahorse among marine animal emojis, starfish, platypus).
``Implausible nonexistent'' items are real entities with weak category association (\eg barnacle, plankton, sea slug for marine emojis).
This distinction tests whether semantic plausibility drives false positive rates.

\subsection{Models}

We test two models from the GPT-4.1 family:

\begin{itemize}[leftmargin=*,itemsep=0pt,topsep=0pt]
    \item \gptfull---OpenAI's full-size model, queried via the OpenAI API.
    \item \gptmini---OpenAI's smaller, faster variant.
\end{itemize}

All queries use temperature 0.0 for deterministic outputs, with a maximum of 300~tokens per response (1,000 for listing experiments).
All API responses are cached to SHA-256-keyed files for reproducibility.

\subsection{Experimental Design}

We conduct five experiments, each targeting a specific aspect of the category-completion hypothesis.

\para{Experiment 1: Direct existence probing.}
We query each model about all 126~emoji items using the direct prompt: ``Does a \{item\} emoji exist in Unicode?''
This measures baseline \fpr across plausibility levels.

\para{Experiment 2: Category density effect.}
We analyze Experiment~1 data grouped by category size to test whether denser categories (\ie those with more real members) produce higher false positive rates.

\para{Experiment 3: Prompt framing.}
We test 20~key items (10~plausible nonexistent across categories) with four prompt framings:
\begin{enumerate}[leftmargin=*,itemsep=0pt,topsep=0pt]
    \item \emph{Direct}: ``Does a \{item\} emoji exist?''
    \item \emph{Indirect}: ``Show me the \{item\} emoji.''
    \item \emph{Confidence}: Direct prompt plus ``Rate your confidence 0--100.''
    \item \emph{Adversarial}: ``I know the \{item\} emoji exists. Confirm its Unicode codepoint.''
\end{enumerate}

\para{Experiment 4: Cross-domain generalization.}
We test plausible-nonexistent items in non-emoji domains (country capitals and chemical elements) to determine whether category-completion false memories extend beyond emojis.

\para{Experiment 5: Category listing.}
We ask models to list all emojis in each category, checking for spontaneous false inclusions. This tests whether generative listing engages different processes than verification.

\subsection{Evaluation Metrics}

We report four metrics grounded in signal detection theory \cite{green1966signal}:

\begin{itemize}[leftmargin=*,itemsep=0pt,topsep=0pt]
    \item \textbf{False positive rate (FPR)}: $P(\text{claims exist} \mid \text{does not exist})$.
    \item \textbf{True positive rate (TPR)}: $P(\text{claims exist} \mid \text{exists})$.
    \item \textbf{$d'$ (d-prime)}: Sensitivity, measuring the separation between hit and false alarm distributions.
    \item \textbf{Criterion ($c$)}: Response bias, where $c > 0$ indicates conservative bias (tendency to say ``No'').
\end{itemize}

\para{Statistical tests.}
We use Fisher's exact test for comparing \fpr between plausibility conditions (small cell counts), chi-square tests for prompt framing comparisons, and Spearman rank correlation for the density--\fpr relationship.
We set $\alpha = 0.05$ as the significance threshold.

\para{Response parsing.}
We classify model responses using a rule-based parser that detects Yes/No indicators in the first line, with fallback heuristics for ambiguous responses. Manual verification of a random sample confirmed $>$95\% parsing accuracy.
